% Part: incompleteness
% Chapter: arithmetization-syntax
% Section: coding-terms

\documentclass[../../../include/open-logic-section]{subfiles}

\begin{document}

\olfileid[fr]{inc}{art}{trm}
\olsection{Codage des termes}

\begin{explain}
Un terme est simplement une suite de symboles d'un certain type :
il se construit inductivement à partir des constantes et des
variables, selon les règles de formation des termes. Puisque nous
pouvons coder les suites de symboles par des nombres, au moyen d'un
codage des symboles et d'un codage des suites de nombres, il n'est
pas difficile d'attribuer des numéros de G\"odel aux termes. La
difficulté est plutôt de montrer que la propriété d'être le numéro
de G\"odel d'un terme bien formé est calculable, et même récursive
primitive.
\end{explain}

!!^{variable}s et !!{constant}s sont les termes les plus simples,
et il est facile de tester si $x$ est le numéro de G\"odel d'un tel
terme : $\fn{Var}(x)$ est vraie si $x$ vaut $\Gn{\Obj v_i}$ pour un
certain~$i$. Autrement dit, $x$ est le code d'une suite de longueur~$1$
dont l'unique élément $(x)_0$ est le code d'!!a{variable}~$\Obj v_i$;
ainsi, $x$ vaut $\tuple{\tuple{1, i}}$ pour un certain~$i$. De même,
$\fn{Const}(x)$ est vraie si $x$ vaut $\Gn{\Obj c_i}$ pour un
certain~$i$. Ces deux relations sont récursives primitives, car, si
un tel $i$ existe, on a nécessairement $i<x$ :
\begin{align*}
  \fn{Var}(x) & \defiff \bexists{i<x}{x = \tuple{\tuple{1, i}}}\\
  \fn{Const}(x) & \defiff \bexists{i<x}{x = \tuple{\tuple{2, i}}}
\end{align*}

\begin{prop}
\ollabel{prop:term-primrec}
Les relations $\fn{Term}(x)$ et $\fn{ClTerm}(x)$, qui sont vraies si
et seulement si $x$ est respectivement le numéro de G\"odel d'un
terme ou d'un terme clos, sont récursives primitives.
\end{prop}

\begin{proof}
Une suite de symboles~$s$ est un terme si et seulement s'il existe
une suite de termes $s_0$, \dots, $s_{k-1} = s$ qui consigne la
formation de~$s$ à partir des !!{constant}s et des !!{variable}s,
selon les règles de formation des termes. Pour que cette suite soit
effectivement une suite de formation, il faut que, pour chaque $i<k$,
on ait l'un des cas suivants :
\begin{enumerate}
\item $s_i$ est !!a{variable}~$\Obj v_j$;
\item $s_i$ est !!a{constant}~$\Obj c_j$;
\item $s_i$ est construit à partir de $n$ termes $t_1$, \dots, $t_n$
  figurant avant la place~$i$, à l'aide d'!!a{function}~$\Obj f^n_j$
  d'arité~$n$.
\end{enumerate}
Pour montrer que la relation correspondante sur les numéros de
G\"odel est récursive primitive, nous devons exprimer cette
condition de façon récursive primitive, c'est-à-dire au moyen de
fonctions et de relations récursives primitives et de
quantifications bornées.

Supposons que $y$ soit le nombre codant la suite $s_0$, \dots,
$s_{k-1}$, c'est-à-dire que
$y = \tuple{\Gn{s_0}, \dots, \Gn{s_{k-1}}}$. Il code une suite de
formation du terme de numéro de G\"odel~$x$ si et seulement si,
pour tout $i<k$, l'un des cas suivants se présente :
\begin{enumerate}
\item $\fn{Var}((y)_i)$;
\item $\fn{Const}((y)_i)$;
\item il existe des indices $n>0$ et $j$, ainsi qu'un nombre
  $z = \tuple{z_1, \dots, z_n}$, tels que chaque $z_l$ soit égal à un
  $(y)_{i'}$ pour un certain $i'<i$ et que
\[
(y)_i = \Gn{\Obj f^n_j(} \concat \fn{flatten}(z) \concat \Gn{)},
\]
\end{enumerate}
et, de plus, $(y)_{k-1} = x$. (La fonction $\fn{flatten}(z)$
transforme la suite $\tuple{\Gn{t_1}, \dots, \Gn{t_n}}$ en
$\Gn{t_1, \dots,
  t_n}$ et est récursive primitive.)

Les indices $j$ et $n$, ainsi que les numéros de G\"odel $z_l$ des termes
$t_l$, sont inférieurs à~$y$. Le code~$z$ de leur suite peut être borné par
$z<\fn{sequenceBound}(y,y)$.\footnote{Précision éditoriale : la source affirme sans justification la borne plus forte $z<y$. La borne récursive primitive indiquée suffit : la suite a $n<y$ éléments, chacun inférieur à~$y$. La variable $j$, omise dans la quantification source, est aussi explicitement liée dans le cas~(3).} Nous pouvons
remplacer $k$ ci-dessus par $\len{y}$. La relation « $y$ est le code
d'une suite de formation du terme de numéro de G\"odel~$x$ » peut
donc être exprimée d'une manière qui montre qu'elle est récursive
primitive.

Il nous reste à nous assurer qu'il existe une borne récursive
primitive pour~$y$. Si $x$ est le numéro de G\"odel d'un terme,
celui-ci possède une suite de formation d'au plus $\len{x}$ termes :
chaque terme de cette suite doit commencer à une certaine place
dans~$s$, et deux sous-termes ne peuvent commencer à la même place.
Le numéro de G\"odel de chaque sous-terme de~$s$ est bien sûr
$\le x$. Il existe donc toujours une suite de formation dont le code
est $\le p_{k-1}^{k(x+1)}$, où $k=\len{x}$.

Pour $\fn{ClTerm}$, il suffit d'omettre le cas des !!{variable}s.
\end{proof}

\begin{prob}
Montrer que la fonction $\fn{flatten}(z)$, qui transforme la suite
$\tuple{\Gn{t_1}, \dots, \Gn{t_n}}$ en $\Gn{t_1, \dots, t_n}$, est
récursive primitive.
\end{prob}

\begin{prop}\ollabel{prop:num-primrec}
  La fonction $\fn{num}(n) = \Gn{\num{n}}$ est récursive primitive.
\end{prop}

\begin{proof}
  Nous définissons $\fn{num}(n)$ par récurrence primitive :
  \begin{align*}
    \fn{num}(0) & = \Gn{\Obj 0}\\
    \fn{num}(n+1) & = \Gn{\prime(} \concat \fn{num}(n) \concat \Gn{)}.
  \end{align*}
\end{proof}

\end{document}
